%! Author = sriadityavedantam
%! Date = 3/30/26

\section{Division}

\begin{theorem}[Division Algorithm in Polynomial Fields]{
    Suppose $\f$ is a field and $a(x),b(x)\in\f[x]$ with $b(x)\neq 0$.
    Then there exist unique $q(x),r(x)\in\f[x]$ such that
    \[
        a(x)=q(x)b(x)+r(x)
    \]
    with $\deg(r(x))<\deg(b(x))$ or $r(x)=0$.
}
    \textbf{Case 1:}
    If $a(x)=0$ or $\deg(a(x))<\deg(b(x))$, then let
    \[
        q(x)=0,\qquad r(x)=a(x).
    \]
    Then
    \[
        a(x)=b(x)\cdot 0+a(x),
    \]
    so the conclusion holds.
    \textbf{Case 2:}
    Suppose $a(x)\neq 0$ and $\deg(a(x))\geq \deg(b(x))$.
    Write
    \[
        a(x)=a_{n}x^n+\hdots,
        \qquad
        b(x)=b_{m}x^m+\hdots
    \]
    with $a_n,b_m\neq 0$ and $n\geq m$.
    Since $\f$ is a field, $b_m^{-1}$ exists.
    Let
    \[
        h(x)\coloneqq a(x)-a_{n}b_m^{-1}x^{n-m}b(x).
    \]
    Then the leading terms cancel, so
    \[
        \deg(h(x))<\deg(a(x)).
    \]
    By strong induction on $\deg(a(x))$, there exist $q_1(x),r(x)\in\f[x]$ such that
    \[
        h(x)=q_1(x)b(x)+r(x)
    \]
    with $\deg(r(x))<\deg(b(x))$ or $r(x)=0$.
    Therefore
    \begin{align*}
        a(x)
        &=a_{n}b_m^{-1}x^{n-m}b(x)+h(x)\\
        &=a_{n}b_m^{-1}x^{n-m}b(x)+q_1(x)b(x)+r(x)\\
        &=\bigl(a_{n}b_m^{-1}x^{n-m}+q_1(x)\bigr)b(x)+r(x).
    \end{align*}
    So if we let
    \[
        q(x)\coloneqq a_{n}b_m^{-1}x^{n-m}+q_1(x),
    \]
    then
    \[
        a(x)=q(x)b(x)+r(x)
    \]
    with $\deg(r(x))<\deg(b(x))$ or $r(x)=0$.

    \textbf{Proof of Uniqueness:}
    Suppose
    \begin{align*}
        a(x)&=q_1(x)b(x)+r_1(x)\\
        &=q_2(x)b(x)+r_2(x),
    \end{align*}
    where $\deg(r_1(x)),\deg(r_2(x))<\deg(b(x))$ or $r_1(x)=r_2(x)=0$.
    Then
    \begin{align*}
        [q_1(x)-q_2(x)]b(x)+[r_1(x)-r_2(x)]&=0,\\
    [q_1(x)-q_2(x)]b(x)&=r_2(x)-r_1(x).
    \end{align*}

    Now either $r_2(x)-r_1(x)=0$, or
    \[
        \deg(r_2(x)-r_1(x))<\deg(b(x)).
    \]
    If $(q_1(x)-q_2(x))b(x)\neq 0$, then since $\f[x]$ is an integral domain,
    \[
        \deg\bigl((q_1(x)-q_2(x))b(x)\bigr)=\deg(q_1(x)-q_2(x))+\deg(b(x))\geq \deg(b(x)).
    \]
    This is impossible, since the right-hand side has degree strictly less than $\deg(b(x))$ or is $0$.
    Therefore
    \[
        (q_1(x)-q_2(x))b(x)=0.
    \]
    Hence $q_1(x)=q_2(x)$.
    Then
    \[
        r_2(x)-r_1(x)=0,
    \]
    so $r_1(x)=r_2(x)$.
\end{theorem}

The Division Algorithm for polynomial fields is a fundamental concept that allows you to divide one polynomial by another, similar to the division algorithm with integers.
This algorithm helps you express one polynomial as a quotient of another polynomial plus a remainder.

\begin{definition}[Logical Divide of Polynomial Fields]
    Let $a(x),b(x)\in\f[x]$ and $b(x)\neq 0$.
    We say $b(x)\mid a(x)$ if there exists a $q(x)\in\f[x]$ such that
    \[
        a(x)=q(x)b(x).
    \]
\end{definition}

\begin{definition}[GCD of Polynomial Fields]
    Suppose $a(x),b(x)\in\f[x]$ are not both $0$.
    We say $d(x)=\gcd(a(x),b(x))$ if $d(x)\mid a(x)$, $d(x)\mid b(x)$, and whenever $c(x)\in\f[x]$ satisfies $c(x)\mid a(x)$ and $c(x)\mid b(x)$, then $c(x)\mid d(x)$.
\end{definition}

Suppose we are in $\Q[x]$.
Let
\[
    a(x)=(x-1)^2
\]
and
\[
    b(x)=(x-1)(x-2).
\]
Then $\gcd(a(x),b(x))=x-1$.
But wait, does not $2x-2\mid a(x)$ and $b(x)$ as well.
We have a problem on our hands.
We have to figure out how to circumvent this solution, and before we can do that, let us go ahead and introduce a new term.

\begin{definition}[Monic]
    If $d(x)\in\f[x]$ has leading coefficient $1$, then $d(x)$ is monic.
\end{definition}

In algebra, monic polynomials are commonly used in the context of irreducible polynomials.
Monic irreducible polynomials have leading coefficient $1$, and this condition simplifies discussions of unique factorization.

\begin{definition}[Polynomial Associates]
    If $c(x),d(x)\in\f[x]$ and
    \[
        c(x)=\beta d(x)
    \]
    for some $\beta\in\f$ with $\beta\neq 0$, we say $c(x)$ and $d(x)$ are associates.
\end{definition}

We can think of associates as polynomial constant multiples.

\begin{theorem}[GCD Theorem]{
    Suppose $a(x),b(x)\in\f[x]$ are not both $0$.
    Let
    \[
        S\coloneqq \{u(x)a(x)+v(x)b(x)\neq 0 : u(x),v(x)\in\f[x]\}.
    \]
    Then there exist $u(x),v(x)\in\f[x]$ such that
    \[
        d(x)=u(x)a(x)+v(x)b(x)
    \]
    and $d(x)=\gcd(a(x),b(x))$.
    Moreover, $S$ has a unique monic polynomial of smallest degree, and this polynomial is $\gcd(a(x),b(x))$.
}
    This theorem also answers the question to our gcd question, which shows that we want to have a monic polynomial of smallest degree as our $\gcd(a(x),b(x))$.
    The set of degrees is a subset of $\z^{\geq 0}$, so let $d(x)$ be a monic polynomial of minimal degree in $S$.
    We need to show that $d(x)\mid a(x)$.
    Using the division algorithm, write
    \[
        a(x)=q(x)d(x)+r(x),
    \]
    where $r(x)=0$ or $\deg(r(x))<\deg(d(x))$.
    Now since $d(x)\in S$, there exist $u(x),v(x)\in\f[x]$ such that
    \[
        d(x)=u(x)a(x)+v(x)b(x).
    \]
    So
    \begin{align*}
        r(x)
        &=a(x)-q(x)d(x)\\
        &=a(x)-q(x)\bigl(u(x)a(x)+v(x)b(x)\bigr)\\
        &=\bigl(1-q(x)u(x)\bigr)a(x)-q(x)v(x)b(x).
    \end{align*}
    Thus if $r(x)\neq 0$, then $r(x)\in S$.
    But then
    \[
        \deg(r(x))<\deg(d(x)),
    \]
    contradicting the fact that $d(x)$ has the least degree among nonzero elements of $S$.
    So $r(x)=0$, and hence $d(x)\mid a(x)$.
    Similarly, $d(x)\mid b(x)$.
    Now suppose $c(x)\mid a(x)$ and $c(x)\mid b(x)$.
    Then $c(x)$ divides every linear combination of $a(x)$ and $b(x)$.
    In particular,
    \[
        c(x)\mid u(x)a(x)+v(x)b(x)=d(x).
    \]
    Therefore $d(x)=\gcd(a(x),b(x))$.
    Finally, gcds are unique up to associates, and among all associates there is exactly one monic polynomial.
    Hence the monic gcd is unique.
\end{theorem}

This theorem also answers the question to our gcd question, which shows that we want to have a monic polynomial of smallest degree as our $\gcd(a(x),b(x))$.
The GCD Theorem for polynomial fields is a fundamental result in abstract algebra that addresses the existence and uniqueness of the greatest common divisor of two polynomials in a polynomial ring over a field.
The theorem establishes a clear and precise method for finding the gcd of polynomials and its properties.

\begin{definition}[Relatively Prime]
    Suppose $a(x),b(x)\in\f[x]$ are not both $0$.
    We say $a(x)$ and $b(x)$ are relatively prime if
    \[
        \gcd(a(x),b(x))=1.
    \]
\end{definition}

\begin{corollary}[Consequence of GCD Theorem]{
    Suppose $a(x),b(x)\in\f[x]$ are relatively prime and $c(x)\in\f[x]$.
    If $a(x)\mid b(x)c(x)$, then $a(x)\mid c(x)$.
}
    By the gcd theorem, we have
    \[
        1=u(x)a(x)+v(x)b(x)
    \]
    for some $u(x),v(x)\in\f[x]$.
    Thus
    \[
        c(x)=c(x)u(x)a(x)+c(x)v(x)b(x).
    \]
    Since $a(x)\mid c(x)u(x)a(x)$ and $a(x)\mid c(x)v(x)b(x)$, it follows that
    \[
        a(x)\mid c(x).
    \]
\end{corollary}

If we let $\ri=\f[x]$, we notice that it has very similar properties to $\z$, such that it has the division and gcd algorithm.
In fact, it also will have relatively prime and an equivalence to primes, but for polynomials.
Let us look at this equivalence.

\begin{definition}[Irreducible]
    A polynomial $p(x)\in\f[x]$ is irreducible if whenever
    \[
        p(x)=a(x)b(x)
    \]
    for $a(x),b(x)\in\f[x]$, then $a(x)$ is an associate of $p(x)$ or $a(x)$ is a unit.
\end{definition}